\chapter{Özet}

Bu proje, bilgisayar bilimi alanındaki Wikipedia makaleleri arasında
makaleden makaleye öneri üreten grafik tabanlı bir öneri sistemini
incelemektedir. Çalışmanın temel amacı, bir kaynak makale verildiğinde onunla
anlamsal ve yapısal olarak ilişkili makaleleri sıralayabilen, daha sonra
istenirse bu sonuçları öğrenme yolu benzeri bir düzende sunabilen bir sistem
oluşturmaktır. Çalışma, kullanıcı profillemesine dayalı klasik öneri
sistemlerinden farklı olarak tamamen sorgu tabanlıdır ve açık Wikipedia verisi
üzerinde çalışır.

Projenin deneysel omurgasını Custom-WikiCS veri kümesi oluşturmaktadır. Bu veri
kümesi, Wiki-CS grafik yapısının gerçek Wikipedia makale metinleri ile
zenginleştirilmiş bir sürümüdür. İlk aşamada 11,701 düğümlü grafik
hazırlanmış, izole düğümler çıkarıldıktan sonra 11,367 etkin düğüm ve 291,039
yönlü kenardan oluşan temizlenmiş grafik elde edilmiştir. Yapısal analizler,
ağın ağır kuyruklu derece dağılımına sahip olduğunu, bağlantılı makalelerin
rastgele makale çiftlerine göre daha yüksek anlamsal benzerlik taşıdığını ve
Louvain ile Infomap gibi yöntemlerin anlamlı topluluklar çıkarabildiğini
göstermiştir. Bu sonuçlar, Wikipedia hiperbağlantı yapısının öneri üretimi için
yararlı bir sinyal taşıdığını doğrulamıştır.

Yöntem tarafında üç ana bileşen kullanılmıştır. İlk olarak, makale metinleri
BAAI/bge-m3 ile sayısal gömülere dönüştürülmüştür. İkinci olarak, grafik
yapısından Node2Vec tabanlı yapısal gömüler elde edilmiştir. Üçüncü olarak da
GCN ve GraphSAGE gibi grafik sinir ağı tabanlı yapısal temel modeller
eğitilmiştir. Çalışmanın erken aşamalarında anlamsal ve yapısal temsil
vektörleri birleştirilerek öneri üretilmiş, daha sonra çoklu model
karşılaştırmaları yapılmıştır. Böylece yalnızca semantik, yalnızca yapısal,
hibrit, birleştirilmiş ve grafik sinir ağı tabanlı yaklaşımlar birlikte
değerlendirilmiştir.

Projenin önemli katkılarından biri, değerlendirme protokolünün daha sonra
yeniden tasarlanmış olmasıdır. İlk değerlendirme düzeninde yönlü bir kenar
ayırma yaklaşımı kullanılırken, yapısal modeller fiilen yönsüz yapı üzerinde
çalışıyordu. Bu durum özellikle Node2Vec sonuçlarının yorumunu etkileyebilecek
metodolojik bir uyumsuzluk oluşturuyordu. Bu nedenle yönsüz yapıya uygun yeni
bir değerlendirme protokolü kurulmuş, 216,123 yönsüz kenar 172,900 eğitim ve
43,223 test kenarı olarak ayrılmış, ek izole düğüm oluşmaması sağlanmış ve
bütün ana yapısal modeller bu düzeltilmiş düzen içinde tekrar
değerlendirilmiştir.

Sonuçlar iki farklı bakış açısını ortaya koymaktadır. İlk değerlendirme
düzeninde düzeltilmemiş Node2Vec çok güçlü görünmekteydi; örneğin Report-6
sonuçlarında Hit@10 değeri \%16.06'ya kadar çıkmıştır. Ancak düzeltilmiş
değerlendirme protokolünde bu üstünlük kaybolmuş, en iyi genel öneri ailesi
düzeltilmiş hibrit modeller olmuştur. Özellikle \textit{BGE pooled + Node2Vec (corrected)} modeli
Hit@10 açısından en güçlü sonuçlardan birini verirken, \textit{BGE-M3 +
Node2Vec (corrected)} modeli MRR, Recall ve MAP gibi ikincil sıralama
ölçütlerinde çok rekabetçi kalmıştır. Buna karşılık ROC-AUC ve PR-AUC gibi bağ
ayırımcılığı ölçütlerinde saf yapısal modellerin hâlâ çok güçlü olduğu
görülmüştür. Bu durum, en iyi öneri modeli ile en iyi küresel bağ ayırma
modelinin aynı şey olmadığını göstermektedir.

Çalışmada sadece sıralama doğruluğu değil, öneri kalitesi de incelenmiştir.
Popülerlik artışı, Gini katsayısı, toplu çeşitlilik, yenilik ve liste içi
uzaklık gibi ölçütler; bazı modellerin doğruluk açısından yakın sonuçlar verse
de çok farklı öneri davranışlarına sahip olabildiğini göstermiştir. Örneğin
GCN Only modeli üst sıralarda zayıf olsa da çeşitlilik ve popülerlik yanlılığı
açısından güçlü bir profil sergilemiştir. BGE-M3 Only ise daha yüksek
popülerlik eğilimi göstermiştir. Düzeltilmiş hibrit modeller ise doğruluk ile
öneri kalitesi arasında en dengeli aile olarak öne çıkmıştır.

Projenin kullanıcıya dönük tarafında yerel bir web arayüzü de
geliştirilmiştir. Bu arayüz başlık arama, model seçimi, ilk 20 önerinin
listelenmesi, makale ön izleme paneli ve isteğe bağlı öğrenme yolu görünümü
sunmaktadır. Öğrenme yolu katmanı başlangıçta küçük bir yerel model ve Ollama
ile denenmiş, ancak tek seferde kararlı yapılandırılmış çıktı üretmekte
yetersiz kaldığı görülmüştür. Bu nedenle mevcut uygulamada sıralama ve erişim
yerel kalırken, öğrenme yolu organizasyonu Gemini tabanlı yapılandırılmış
çıktı, uygulama tarafında katı doğrulama ve başarısızlık durumunda sezgisel
geri dönüş mekanizması ile güncellenmiştir. Böylece LLM katmanı, asıl öneri
mantığından ayrı, sonradan eklenen bir sunum katmanı olarak ele alınmıştır.

Bu çalışma; açık veri, yerel önbellekleme ve tüketici sınıfı donanım ile
yürütülmüş olması bakımından ekonomik ve tekrarlanabilir bir deneysel çerçeve
sunmaktadır. Kişisel veri kullanılmaması, kullanıcı takibi yapılmaması ve
önerilerin önceden hesaplanmış benzerlik uzayları üzerinden üretilmesi etik ve
güvenlik açısından önemli avantajlar sağlamaktadır. Sonuç olarak proje,
Wikipedia makaleleri için grafik destekli öneri üretimi konusunda hem
metodolojik hem uygulamalı bir katkı ortaya koymakta; aynı zamanda
değerlendirme tasarımının sonuç yorumunu nasıl değiştirebildiğini açık biçimde
göstermektedir.
